
\documentclass[
% twocolumn,
% hf,
]{./CEURART/ceurart}

%%
%% One can fix some overfulls
\sloppy

%%
%% Minted listings support 
%% Need pygment <http://pygments.org/> <http://pypi.python.org/pypi/Pygments>
\usepackage{listings}
%% auto break lines
\lstset{breaklines=true}

\usepackage{lipsum}
%%
%% end of the preamble, start of the body of the document source.
\begin{document}

%%
%% Rights management information.
%% CC-BY is default license.
\copyrightyear{2026}
\copyrightclause{Copyright for this paper by its authors.
	Use permitted under Creative Commons License Attribution 4.0
	International (CC BY 4.0).}

%%
%% This command is for the conference information
\conference{NORA'26: Workshop on KNOwledge GRaphs & Agentic Systems Interplay}

%%
%% The "title" command
\title{Open Data Fusion for Approximate True Cost Accounting using Knowledge Graphs}

% \tnotemark[1]
% \tnotetext[1]{You can use this document as the template for preparing your
% 	publication. We recommend using the latest version of the ceurart style.}

%%
%% The "author" command and its associated commands are used to define
%% the authors and their affiliations.
\author[1]{Bileam Scheuvens}[
	orcid=0009-0006-5632-9421,
	email=bileam.scheuvens@student.uni-tuebingen.de,
	% url=https://yamadharma.github.io/,
]
\cormark[1]
% \fnmark[1]
\address[1]{Eberhard Karls Universität Tübingen, Geschwister-Scholl-Platz, 72074 Tübingen, Germany}
\cortext[1]{Corresponding author.}
% \address[2]{Joint Institute for Nuclear Research,
% 6 Joliot-Curie, Dubna, Moscow region, 141980, Russian Federation}

\author[1]{Harry Scells}[
	orcid=0000-0001-9578-7157
	email=harrisen.scells@uni.tuebingen.de,
	url=https://scells.me/,
	% TODO title?
]
% \fnmark[1]
% \address[1]{Eberhard Karls Universität Tübingen, Geschwister-Scholl-Platz, 72074 Tübingen, Germany}

%
%% Footnotes
% \cortext[1]{Corresponding author.}
% \fntext[1]{These authors contributed equally.}

%%
%% The abstract is a short summary of the work to be presented in the
%% article.
\begin{abstract}
	A clear and well-documented \LaTeX{} document is presented as an
	article formatted for publication by CEUR-WS in a conference
	proceedings. Based on the ``ceurart'' document class, this article
	presents and explains many of the common variations, as well as many
	of the formatting elements an author may use in the preparation of
	the documentation of their work.
\end{abstract}

%%
%% Keywords. The author(s) should pick words that accurately describe
%% the work being presented. Separate the keywords with commas.
\begin{keywords}
	QLever \sep
	Data Fusion \sep
	Knowledge Graphs \sep
	True Cost Accounting
\end{keywords}

%%
%% This command processes the author and affiliation and title

%% information and builds the first part of the formatted document.

\maketitle

\section{Introduction}

\subsection{Related Work}

\section{Methodology}


\subsection{Conceptual Workflow}

\subsection{Implementation}


\begin{figure}
	\centering
	\includegraphics[width=\linewidth-150pt]{architecture-1.png}
	\caption{Architecture Diagram}
\end{figure}


\section{Discussion}

\subsection{Limitations}

\subsection{Future Work}


%%
%% The acknowledgments section is defined using the "acknowledgments" environment
%% (and NOT an unnumbered section). This ensures the proper
%% identification of the section in the article metadata, and the
%% consistent spelling of the heading.
\begin{acknowledgments}
	The authors gratefully acknowledge the computing time granted by the KISSKI project. Some calculations for this research were conducted with computing resources under the project HiddenCostReport.
\end{acknowledgments}

%% The declaration on generative AI comes in effect
%% in Janary 2025. See also
%% https://ceur-ws.org/GenAI/Policy.html
\section*{Declaration on Generative AI}
The authors have not employed any Generative AI tools.
\newline
%%
%% Define the bibliography file to be used
\bibliography{paper}
% \bibliography{../report/hiddencostreport.bib}

%%
%% If your work has an appendix, this is the place to put it.
\appendix

\section{Online Resources}


The sources for the ceur-art style are available via
\begin{itemize}
	\item \href{https://github.com/yamadharma/ceurart}{GitHub},
	      % \item \href{https://www.overleaf.com/project/5e76702c4acae70001d3bc87}{Overleaf},
	\item
	      \href{https://www.overleaf.com/latex/templates/template-for-submissions-to-ceur-workshop-proceedings-ceur-ws-dot-org/pkfscdkgkhcq}{Overleaf
		      template}.
\end{itemize}

\end{document}

%%
%% End of file

